% ID: SEL-D0040-OB-E01
% ESTADO: cotejada. Segunda pasada visual y borrador compilado revisados.
% Fuente: Selectividad/matematicas_aplicadas/Ua5MATEMÁTICAS APLICADAS examen.pdf
% SHA-256: ac80a164e049669cd01083ca5adaddfa90efbf556b050de2d99f294ab55670c9
% Página: 2; opción B; ejercicio original 1.

Sea el siguiente conjunto de inecuaciones:
\[
x-3y\leq8;\quad3x+2y\geq15;\quad x+3y\leq12;\quad x\geq0;\quad y\geq0.
\]
\begin{enumerate}[label=\alph*)]
\item \textbf{(1 punto)} Dibuje el recinto del plano determinado por estas inecuaciones.
\item \textbf{(1 punto)} Determine los vértices de este recinto.
\item \textbf{(0.5 puntos)} Maximice la función $F(x,y)=5x+9y$ en este recinto, indicando el punto o puntos donde se alcanza ese máximo.
\end{enumerate}
